\section{Small transport with nontrivial holonomy}\label{sec:transport}

Retain the finite flag complex $L$, the homomorphism
$\alpha:\pi_1(L)\to H=\SU(2)$, and the level complex $X$ of
Proposition~\ref{prop:model}.  Write $V$ for the vertex set of $L$ and
$a_v$ for the corresponding standard generator of $P=A_L$.
We realize $|L|$ in $\R^V$ by barycentric coordinates and give it the
metric induced by the $\ell^1$ norm.  Our aim is the following.

\begin{proposition}\label{prop:small-labels}
For every $\epsilon>0$, the oriented edges of $X$ admit labels
$h_e\in H$ such that
\[
 h_{\bar e}=h_e^{-1},\qquad |h_e-1|<\epsilon,
\]
the product around every triangular two-cell is $1$, and the product
around some closed edge path is different from $1$.
\end{proposition}

For an abelian target, products along closed paths would factor through
$H_1(X;\Z)=0$ and would all be trivial.  The nonabelian target is
therefore essential.

The construction has two parts.  A direction map on reduced words of
$P$ changes slowly far from the identity.  Translating the level $X$
far from the identity therefore gives a map of its graph to $L$ with
small edge images.  Its image nevertheless contains prescribed loops
of $L$, allowing $\alpha$ to detect nontrivial holonomy.

The long commuting-power paths below lie along expanded copies of $L$
in level sets, formed from commuting coordinate directions as in
Bestvina and Brady
\cite[Definition~8.8 and the following paragraph]{BB1997}.

\subsection{Reduced traces and their direction vectors}

Two signed standard generators are called \emph{independent} when
their underlying vertices are distinct and adjacent in $L$.
In particular, two occurrences of the same generator are dependent,
regardless of their signs.  A \emph{trace} is a word modulo exchanges
of adjacent independent letters, as in the theory of partially
commutative monoids \cite[Chapter~I, Section~2]{CartierFoata1969}.  Occurrences can be followed through
these exchanges.  The \emph{occurrence poset} of a word orders every
dependent pair in its order of appearance and takes the transitive
closure.  It records the trace: its linear extensions are precisely
the representatives of that trace; see also Krattenthaler's appendix
to the 2006 reedition of \cite[Definition~2.2 and Section~3 of the appendix]{CartierFoata1969}.
Indeed, exchanges preserve this
poset, and any two linear extensions of a finite poset are related by
exchanges of adjacent incomparable elements.  Here incomparable
occurrences have independent directions.

The following signed-letter version of the graph-product normal form
\cite[Theorem~3.9]{Green1990} records individual occurrences rather than
vertex-group syllables. We include its reduction proof because that
occurrence information is used in the direction-vector estimate.

\begin{lemma}\label{lem:trace-normal-form}
Every element of $P$ has a unique reduced trace, obtained by deleting
inverse pairs that can be made adjacent by exchanges.  Its length is
the word length in the standard generators.  Right multiplication by
a signed standard generator changes this occurrence poset by either
adjoining a maximal occurrence or deleting a maximal occurrence.
\end{lemma}

\begin{proof}
An inverse pair can be deleted exactly when every occurrence between
its endpoints is independent of its direction.  This criterion is
independent of the representative: the order of an occurrence
dependent on that direction relative to either endpoint cannot change
under exchanges.  Deletion gives the trace represented by removing the
two occurrences.

We check the local confluence of deletion.  For two disjoint deletable
pairs, either deletion leaves the other available, and both orders
remove exactly the same four occurrences.  If two pairs share an
occurrence, they cannot share their first endpoint or their last
endpoint, because an intervening occurrence of the same direction
would prevent deletion.  The relevant subword consequently has the
form
\[
 x B x^{-1} C x
\]
or the same form with $x$ replaced by $x^{-1}$, where every letter of
$B$ and $C$ is independent of $x$.  The two deletions leave $BCx$ and
$xBC$, which are the same trace.  Since deletion decreases length,
induction on length using these local diamonds proves uniqueness of
the terminal trace.  Exchanges and adjacent inverse insertions or
deletions preserve that terminal trace.  These operations generate
equality in the defining presentation of $P$, so the terminal trace
depends only on the group element.  Every word representing that
element reduces to it, proving the assertion about word length.

Write $|g|_P$ for word length in the signed standard generators.
Let $w$ be reduced of length $n$, and append a signed generator $x$.
If the result is reduced, its new occurrence is maximal.  Otherwise
an available deletion must involve the new last occurrence, since a
deletion involving only old occurrences would already have been
available in $w$.  Its inverse mate $y$ has only independent
occurrences after it in $w$, and is therefore maximal in the old
occurrence poset.  Deleting $y$ and the new $x$ leaves a word of length
$n-1$.  If $g$ is the element represented by $w$, the word-metric
inequality gives
\[
 |gx|_P\geq |g|_P-1=n-1.
\]
Thus the remaining word is reduced.  Deleting a maximal occurrence
also leaves exactly the induced order on the retained occurrences:
no dependency chain between two retained occurrences can pass through
a maximal occurrence.  This proves the last assertion.
\end{proof}

Fix a total order on $V$.  In a nonempty reduced trace, the minimal
occurrences have distinct pairwise independent directions, hence
their directions span a simplex of $L$.  Assign each occurrence $p$
to the first direction, in our fixed order, among the minimal
occurrences below or equal to $p$.  This set is nonempty because the
poset is finite.  Let $c_v(g)$ be the number assigned to $v$, and set
\begin{equation}\label{eq:direction-map}
 b(g)=\frac{1}{|g|_P}\sum_{v\in V}c_v(g)e_v\in |L|,
 \qquad g\ne 1.
\end{equation}
Here $e_v$ is the barycentric vertex corresponding to $v$.
Each minimal occurrence is assigned to its own direction, so the
support of $b(g)$ is exactly the set of minimal directions.

\begin{lemma}\label{lem:trace-direction}
Suppose $g$ and $gx$ are nonidentity, where $x$ is a signed standard
generator.  Their direction vectors belong to a common simplex of
$L$.  If their word lengths are $n$ and $n+1$, in either order, then
\begin{equation}\label{eq:trace-direction-bound}
 \norm{b(gx)-b(g)}_1\leq \frac{2}{n+1}.
\end{equation}
\end{lemma}

\begin{proof}
By Lemma~\ref{lem:trace-normal-form}, one occurrence poset is obtained
from the other by adjoining a maximal occurrence.  Every retained
occurrence has the same principal downset before and after this
operation: a maximal occurrence cannot lie on a dependency chain
ending at a different occurrence.  The globally minimal predecessors
of a retained occurrence are precisely the minimal elements of its
principal downset.  Its assigned direction is therefore unchanged.
This also covers a changed occurrence that is both minimal and
maximal: it is isolated and has no retained successors.

Consequently, if $c$ is the count vector for the shorter trace, the
longer trace has count vector $c+e_v$ for some $v\in V$.  Since
$\norm{c/n}_1=1$,
\[
 \norm{\frac{c+e_v}{n+1}-\frac{c}{n}}_1
 =\frac{\norm{e_v-c/n}_1}{n+1}\leq\frac{2}{n+1}.
\]
For the common-simplex assertion, adjoining a maximal occurrence
preserves all old minimal occurrences.  If the new occurrence is
not minimal, the minimal set is unchanged.  If it is also
minimal, it is isolated, and its direction is independent of every
old occurrence: dependent occurrences would be comparable.  Thus
the union of the two sets of minimal directions is a clique.
Flagness of $L$ completes the proof.
\end{proof}

The direction vector changes by at most the reciprocal word-length
bound in Equation~\eqref{eq:trace-direction-bound} under a generator
step. We next apply this estimate on a translate of the level graph
where all relevant words are long. At the same time, we retain exact
control of certain paths, so that making every edge image small does
not erase the loops detected by $\alpha$.

\subsection{A slowly varying map of the level graph}

Let $N\geq2$ and choose $d\in P$ with $\lambda(d)=N$.  Define a map on
the vertices of $X$ by
\begin{equation}\label{eq:translated-direction-map}
 f_N(g)=b(d^{-1}g),\qquad g\in G.
\end{equation}
It is defined because $\lambda(d^{-1}g)=-N$.  Every edge of $X$ is
the diagonal of a commuting square in $E$ with corner heights
$-1,0,0,1$.  If its endpoints are $g,h$, write $u$ for the square's
unique corner of height $1$.  The two steps $g,u,h$ are standard
generator steps in $E$.  Their translates by $d^{-1}$ have heights
$-N,1-N,-N$, hence word lengths at least $N,N-1,N$.
The larger word length in each step is therefore at least $N$.
Lemma~\ref{lem:trace-direction} gives
\[
 \norm{b(d^{-1}g)-b(d^{-1}u)}_1\leq 2/N,
 \qquad
 \norm{b(d^{-1}u)-b(d^{-1}h)}_1\leq 2/N.
\]
Each pair belongs to a common simplex, so the two straight segments
joining them lie in $|L|$.

Map the edge from $g$ to $h$ to this two-segment path, with the
following convention: whenever $f_N(g)$ and $f_N(h)$ themselves lie
in a common simplex, use their single straight segment instead.
This segment is unambiguous even if several simplices contain the
endpoints: it is the same segment in barycentric coordinates and lies
in their common face.  Choose the reverse path for the reverse
orientation.  These choices define a continuous map
$f_N:X^{(1)}\to |L|$, since they agree at vertices and the domain is a
CW complex.  Every edge image has diameter at most $4/N$.  Moreover,
any two of the three edge images of a triangle meet at their common
vertex image.  Therefore
\begin{equation}\label{eq:small-triangle-image}
 \operatorname{diam} f_N(\partial\sigma)\leq 8/N
 \qquad\text{for every triangular two-cell $\sigma$ of $X$.}
\end{equation}

There is also exact control on selected long paths.  If $\{q,r\}$ is
an edge of $L$, put
\begin{equation}\label{eq:commuting-power-path}
 p_j=d\,a_q^{-(N-j)}a_r^{-j},\qquad 0\leq j\leq N.
\end{equation}
These vertices have height zero.  Consecutive vertices are joined by
an edge of $X$: the square based at $p_j a_r^{-1}$ in the positive
$q,r$ directions has level-zero corners $p_j$ and $p_{j+1}$.
The trace of $d^{-1}p_j$ consists of two mutually independent chains
of negative occurrences, one in each direction, with lengths $N-j$
and $j$; a chain of length zero is omitted.  Every occurrence is
assigned to its own direction.  Hence
\begin{equation}\label{eq:exact-subdivision}
 f_N(p_j)=\left(1-\frac{j}{N}\right)e_q+\frac{j}{N}e_r.
\end{equation}
Our straight-segment convention means that this level path maps
exactly to the edge $[q,r]$, subdivided into $N$ equal parts; see
Figure~\ref{fig:commuting-path}.

\begin{figure}[ht]
 \centering
 \begin{tikzpicture}[x=0.72cm,y=0.72cm,>=Stealth]
  \draw[step=1,gray!35,thin] (0,0) grid (4,4);
  \draw[blue!70!black,very thick] (0,4)--(4,0);
  \foreach \j in {0,...,4}{
   \fill[blue!70!black] (\j,{4-\j}) circle (2.2pt);
  }
  \node[above left] at (0,4) {$p_0=d a_q^{-N}$};
  \node[below right] at (4,0) {$p_N=d a_r^{-N}$};
  \node[below] at (1.6,-0.65) {Level path in a commuting flat};
  \draw[->,thick] (4.6,2.2)--node[above] {$f_N$}(6.3,2.2);
  \draw[blue!70!black,very thick] (7,2)--(11,2);
  \foreach \j in {0,...,4}{
   \fill[blue!70!black] ({7+\j},2) circle (2.2pt);
  }
  \node[below] at (7,1.9) {$q$};
  \node[below] at (11,1.9) {$r$};
  \node[above] at (8,2.15) {$\frac14$};
  \node[above] at (9,2.15) {$\frac12$};
  \node[above] at (10,2.15) {$\frac34$};
  \node[below] at (9,0.65) {The edge $[q,r]\subset L$};
 \end{tikzpicture}
 \caption{The special path~\eqref{eq:commuting-power-path}, shown for
 $N=4$, joins $d a_q^{-N}$ to $d a_r^{-N}$ across level-cut squares.
 Formula~\eqref{eq:exact-subdivision} sends it to the exact subdivision
 of $[q,r]$.  Here $a_q,a_r$ commute and $\lambda(d)=N$.
 Each blue diagonal step multiplies by $a_q a_r^{-1}$, preserving
 height zero; $f_N$ is shown on this level path.}
 \label{fig:commuting-path}
\end{figure}

We have obtained both properties needed from $f_N$: every triangle
boundary has small image, and selected closed paths trace prescribed
loops of $L$ exactly. The representation $\alpha$ will turn these two
properties into flatness on triangles and nontrivial closed-path
holonomy, respectively.

\subsection{Pulling back transport}

\begin{proof}[Proof of Proposition~\ref{prop:small-labels}]
Choose a base vertex $q_0$ for $\alpha$ and a base lift of $q_0$ in
the universal cover $\widetilde L\to L$.  Choose the left deck
action convention so that lifting a loop representing $\gamma$ from
that base lift ends at $\gamma$ times the lift.  There is a
continuous map
\begin{equation}\label{eq:developing-map}
 D:\widetilde L\longrightarrow H,
 \qquad D(\gamma\xi)=\alpha(\gamma)D(\xi).
\end{equation}
To construct it, choose images for representatives of the vertex
orbits and extend equivariantly.  On representatives of the edge
orbits, extend between the prescribed endpoints using path
connectedness of $H$.  Extend over representatives of the two-cell
orbits using $\pi_1(H)=0$.  The deck action is free on cells, so
equivariance introduces no ambiguity; $L$ has dimension two, so there
are no further extensions.  The CW topology gives continuity.

Fix $\epsilon>0$.  There is a finite open cover $(U_i)$ of $L$ such
that each $U_i$ is evenly covered and $D$ has image of diameter less
than $\epsilon$ on each sheet over $U_i$.  Indeed, around any lift
of a point, continuity of $D$ gives a sufficiently small such
neighborhood on one sheet.  All other sheets are its deck translates,
and left multiplication by $\alpha(\gamma)$ preserves quaternion
distance.  Compactness of $L$ then gives a finite subcover.  Let
$\ell>0$ be a Lebesgue number in the diameter form: every subset of
$L$ of diameter less than $\ell$ lies in some $U_i$.  Choose $N\geq2$
so large that $8/N<\ell$, and construct $f_N$ as above.

For an oriented edge $e$ of $X$, lift its assigned path $f_N(e)$ from
$\xi$ to $\xi'$ in $\widetilde L$, and define
\begin{equation}\label{eq:edge-transport}
 h_e=D(\xi)^{-1}D(\xi').
\end{equation}
Changing the lift multiplies both $D$ values on the left by the same
$\alpha(\gamma)$, so the label is independent of the lift.  Reversing
the path inverts the label.  Its image lies in some $U_i$ because its
diameter is at most $4/N<\ell$.  The lift stays in one sheet, and
therefore
\[
 |h_e-1|=|D(\xi')-D(\xi)|<\epsilon.
\]

The whole image of the boundary of a triangular two-cell lies in one
$U_i$ by~\eqref{eq:small-triangle-image}.  The inverse homeomorphism
of a sheet over $U_i$ lifts this boundary to a closed loop in that
sheet.  In multiplying the three edge
labels, choose these consecutive lifts; the intermediate $D$ values
cancel, giving product $1$.

Finally, choose a finite edge loop based at $q_0$,
$q_0,q_1,\ldots,q_k=q_0$ in $L$ whose class $\gamma$ has
$\alpha(\gamma)\ne1$.  Such a loop exists because $\alpha$ is
nontrivial and every element of the fundamental group of a simplicial
complex is represented by a finite edge loop.  For each successive
pair, concatenate the path~\eqref{eq:commuting-power-path}.  The
endpoint of each piece is $d a_{q_{i+1}}^{-N}$, the initial vertex of
the next piece, so this is a closed edge path in $X$ based at
$d a_{q_0}^{-N}$.  By~\eqref{eq:exact-subdivision}, its image under
$f_N$ is precisely the chosen loop with each edge subdivided.
Lift that image from the chosen base lift $\xi$ of $q_0$.  The terminal lift is
$\gamma\xi$, and telescoping~\eqref{eq:edge-transport} gives the
path product
\[
 D(\xi)^{-1}D(\gamma\xi)
 =D(\xi)^{-1}\alpha(\gamma)D(\xi)\ne1.
\]
This proves all assertions of the proposition.
\end{proof}
